\documentclass[a4paper,12pt]{article}
\usepackage[utf8]{inputenc}
\usepackage[russian]{babel}
\usepackage{amsmath, amssymb}
\usepackage{graphicx}
\usepackage{geometry}
\usepackage{booktabs}
\usepackage{array}
\usepackage{indentfirst}

\begin{document}

% --- ТИТУЛЬНЫЙ ЛИСТ ---
\begin{titlepage}
    \centering
    Федеральное государственное автономное образовательное учреждение высшего образования \\
    «Национальный исследовательский Нижегородский государственный университет им. Н.И. Лобачевского» \par
    \vspace{4cm}
    \Large \textbf{Отчёт по лабораторной работе} \\
    \Huge \textbf{«Поляризация»} \par
    \vspace{3cm}
    \large
    \textit{Выполнили:} \\
    Кукушкин Юрий, Калинин Дмитрий, \\
    Можаров Андрей, Зинягин Алексей, \\
    Рассулов Илья \par \vspace{1em}
    \textit{Проверил:} \\
    Старобор Алексей Викторович \par
    \vfill
    2025 г.
\end{titlepage}

% --- РАЗДЕЛ 1: ПОСТОЯННАЯ ВЕРДЕ ---
\section{Постоянная Верде среды}

Постоянной Верде называют коэффициент пропорциональности между углом поворота плоскости поляризации $\theta$ и интегралом от продольной компоненты напряженности магнитного поля $H_z$ по длине кристалла $L$:
\begin{equation}
\theta = V \int_{-L/2+\epsilon_0}^{L/2+\epsilon_0} H_z(z) dz = L \cdot H_z \cdot V
\end{equation}

Для нахождения постоянной проводились вычисления интеграла магнитного поля методом трапеций для различных сред:

\begin{itemize}
    \item \textbf{Оксид тербия ($Tb_2O_3$):} $L = 0.7$ мм, $\int_{-0.35}^{0.35} H_z(z) dz = 9.462862$ Кэ$\cdot$мм.
    \item \textbf{Тербий-алюминиевый гранат ($Tb_3Al_5O_{12}$):} $L = 3$ мм, $\int_{-1.5}^{1.5} H_z(z) dz = 40.509526$ Кэ$\cdot$мм.
    \item \textbf{Тербий-галлиевый гранат ($Tb_3Ga_5O_{12}$):} $L = 9.6$ мм, $\int_{-4.8}^{4.8} H_z(z) dz = 128.220031$ Кэ$\cdot$мм.
\end{itemize}

$$
    \theta = \varphi \hspace{1em} V = \dfrac{\varphi}{\int_{-L/2}^{L/2} H_z(z)dz}
$$

\begin{table}[h!]
\centering
\begin{tabular}{lccc}
\toprule
Материал среды & Длина ($L$), мм & Угол поворота ($\varphi$) & V, град/(Кэ$\cdot$мм\\
\midrule
$Tb_3Ga_5O_{12}$ & 9,6 & $148^{\circ} 40'$ ($391^{\circ} 20'$ полн.) & 3,05 \\
$Tb_3Al_5O_{12}$ & 3 & $174^{\circ} 20'$ & 4,3 \\
$Tb_2O_3$        & 0,7 & $41^{\circ} 40'$ & 4,4 \\
\bottomrule
\end{tabular}
\end{table}

% --- РАЗДЕЛ 2: ГЛЮКОЗА ---
\section{Удельное вращение водного раствора глюкозы}

Плотность глюкозы: $\rho_\text{глюк.} = 1,562 \dfrac{\text{г}}{\text{см}^3}$.
$$
    C' = \dfrac{V_\text{глюкозы}}{V_\text{раствора}}\rho_\text{глюкозы}
$$
$$
    \Delta V = \Delta V_\text{глюкозы} = \Delta V_\text{раствора} = 0,05 \text{ мл}
$$
$$
    \Delta C' = \dfrac{\rho_\text{глюкозы}\Delta V_\text{глюкозы}}{\Delta V_\text{раствора}} \sqrt{1 + \left(\dfrac{V_\text{глюкозы}}{V_\text{раствора}}\right)^2}
$$

Удельное вращение $[\alpha]$ определяется как вращение плоскости поляризации слоем вещества толщиной $L=1$ дм при концентрации $C = 1$ г/мл. Формула расчета:
\begin{equation}
[\alpha] = \frac{\alpha \cdot 100}{L \cdot C}
\end{equation}

Для измерений использовалась установка, включающая лазер, кювету ($L=40$ мм), вращающийся поляризатор и осциллограф.

\subsection*{Результаты измерений для глюкозы}

\begin{table}[h]
\centering
\begin{tabular}{cccccc}
\toprule
Концентрация & $V_{\text{р-ра}}$ (мл) & $V_{\text{воды}}$ (мл) & $V_{\text{глюк}}$ (мл) & $\varphi, ^\circ$ & $C$, г/см$^3$ \\
\midrule
40\% & 5 & 3 & 2 & 19.46 & 0,62 \\
35\% & 5,71 & 3,71 & 2 & 18.00 & 0,54 \\
30\% & 6,6 & 4,6 & 2 & 16.00 & 0.47 \\
25\% & 8 & 6 & 2 & 13.30 & 0,39 \\
20\% & 10 & 8  & 2 & 10.24 & 0,31 \\
15\% & 13,3 & 11,3 & 2 & 7.38  & 0,23 \\
10\% & 20 & 18 & 2 & 6.00 & 0,16 \\
\bottomrule
\end{tabular}
\caption{Результаты измерений для глюкозы с учетом дополнительных данных}
\end{table}

На основе аппроксимации графика $\varphi_{мин}(C)$ получен угловой коэффициент $\langle k \rangle = 31,2 \dfrac{\text{град} \cdot \text{см}^3}{\text{г}}$. 
Итоговое удельное вращение: $[\alpha] = (78 \pm 2) \frac{\text{град} \cdot \text{см}^2}{\text{г}}$.

% --- РАЗДЕЛ 3: КРИСТАЛЛ DKDP ---
\section{Кристалл DKDP: Двулучепреломление}

% Разность фаз между обыкновенной и необыкновенной волнами
\begin{equation}
\varphi = 2\pi l \frac{n' - n''}{\lambda}
\end{equation}

% Итоговое выражение для разности фаз с учетом угла падения
\begin{equation}
\varphi = 2\pi \frac{d}{\lambda \cos \rho} (n_0 - n_e) \sin^2 \rho
\end{equation}

% Разность фаз в приближении малых углов (x << f)
\begin{equation}
\varphi = 2\pi \frac{d(n_0 - n_e)}{\lambda f} \frac{x^2}{\sqrt{f^2 + x^2}}
\end{equation}

% Формула для определения разности показателей преломления
\begin{equation}
(n_0 - n_e) = \frac{\varphi \lambda}{2\pi d} \frac{f^2}{x^2}
\end{equation}

% Условия экстремумов (максимумы и минимумы)
\begin{itemize}
    \item \textbf{Для максимумов ($\varphi = 2\pi k$):} 
    \begin{equation}
    (n_0 - n_e) = \frac{k \lambda}{d} \frac{f^2}{x^2}
    \end{equation}
    
    \item \textbf{Для минимумов ($\varphi = (2k + 1)\pi$):} 
    \begin{equation}
    (n_0 - n_e) = \frac{(2k + 1) \lambda}{2d} \frac{f^2}{x^2}
    \end{equation}
\end{itemize}

% Расчет погрешностей
\begin{equation}
\delta(n_0 - n_e) = \delta d + 2\delta f + 2\delta x
\end{equation}
\begin{equation}
\Delta(n_0 - n_e) = (n_0 - n_e) \cdot \delta(n_0 - n_e)
\end{equation}

Параметры кристалла DKDP: $\lambda = 650$~нм , $d = 3$~мм , $F = 96$~мм, $n_0 = 1,4931$.

$\Delta f = 1$ мм, $\Delta x = 0,05$ мм, $\Delta d = 0,1$ мм.   

\begin{table}[h!]
\centering
\begin{tabular}{ccccc}
\toprule
Номер кольца & $x$, мм & $(n_0 - n_e)$ & $\delta (n_0 - n_e)$ & $\Delta (n_0 - n_e)$ \\
\midrule
1max & 5,9  & 0,057 & 0,071 & 0,004 \\
1min & 8,65 & 0,040 & 0,066 & 0,003 \\
2max & 10,2 & 0,038 & 0,064 & 0,002 \\
2min & 13,13& 0,029 & 0,062 & 0,002 \\
3max & 14,1 & 0,030 & 0,061 & 0,002 \\
3min & 15,2 & 0,030 & 0,061 & 0,002 \\
4max & 15,7 & 0,032 & 0,061 & 0,002 \\
4min & 16,8 & 0,032 & 0,060 & 0,002 \\
5max & 17,8 & 0,032 & 0,060 & 0,002 \\
5min & 18,6 & 0,032 & 0,060 & 0,002 \\
6max & 19,8 & 0,031 & 0,059 & 0,002 \\
6min & 20,15& 0,032 & 0,059 & 0,002 \\
7max & 20,7 & 0,033 & 0,059 & 0,002 \\
7min & 21,5 & 0,032 & 0,059 & 0,002 \\
\midrule
 & & 0,034 & 0,061 & 0,002 \\
\bottomrule
\end{tabular}
\end{table}

Среднее, соответственно, $(n_0 - n_e) = (0,034 \pm 0,002)$.

\end{document}